\section{Aufgaben}

\Title{Multiple Choice}

\Title{Folgen und Reihen}

Sei $(a_n), (b_n)$ zwei Folgen, welche Aussgaen sind korrekt?
\begin{itemize}
    \item[$\square$]
    Wenn $a_n \leq b_n$ für alle $n \in \N$ und $(a_n)$ konvergent ist, so ist auch $(b_n)$ konvergent.
    \item[\rlap{\raisebox{0.3ex}{\hspace{0.4ex}\tiny \ding{52}}}$\square$]
    Wenn $0 \leq b_n \leq a^3_n$ für alle $n \in \N$ und $a_n \to 0$ für $n \to \infty$, dann
    konvergiert auch $(b_n)$.
    \item[$\square$]
    Wenn $|a_n - b_n| \leq \frac{1}{n}$ für alle $n \in N$, dann konvergieren beide Folgen.
    \item[$\square$]
    Wenn $(b^2_n)$ konvergiert, dann konvergiert auch $(b_n)$.
\end{itemize}

Seien $(a_n), (b_n)$ Folgen. Dann gilt:
\begin{itemize}
    \item[\rlap{\raisebox{0.3ex}{\hspace{0.4ex}\tiny \ding{52}}}$\square$]
    Falls $(a_n)$ konvergiert und $(b_n)$ nicht, so ist $(a_n + b_n)$ nicht konvergent.
    \item[$\square$]
    Ist $(a_n)$ beschränkt und $(b_n)$ nicht konvergent, so ist $(a_n \cdot b_n)$ nicht konvergent.
    \item[\rlap{\raisebox{0.3ex}{\hspace{0.4ex}\tiny \ding{52}}}$\square$]
    Falls $(a_n)$ beschränkt und $\Limit b_n = 0$, so folgt, dass $\Limit a_n \cdot b_n = 0$. 
\end{itemize}

Sei $f : \R \to [0, \infty[$, so dass $\lim_{x \to 0}f(x) \neq 0$. Dann gilt:
\begin{itemize}
    \item[$\square$]
    $\forall \epsilon > 0, \exists \delta > 0$, so dass $0 < |x| <\delta \Rightarrow f(x) > \epsilon$.
    \item[\rlap{\raisebox{0.3ex}{\hspace{0.4ex}\tiny \ding{52}}}$\square$]
    Es existiert eine Folge $(x_n)$ mit $\Limit x_n = 0$ und ein $\epsilon$, so dass $|f(x_n)| > \epsilon, n \in \N$.
    \item[\rlap{\raisebox{0.3ex}{\hspace{0.4ex}\tiny \ding{52}}}$\square$]
    Für all $x \in \R$ gilt $f(x) > 0$.
    \item[$\square$]
    Für jede Folge $(x_n)$ mit $\Limit x_n = 0$ gilt $\Limit f(x_n) \neq 0$.
\end{itemize}

Seie $\Reihe a_n$ eine Reihe. Dann gilt:
\begin{itemize}
    \item[$\square$]
    Falls $\forall \epsilon > 0, \exists N \geq 1$, so dass $\sum_{k = n}^{n+100} |a_k| < \epsilon, \forall n \geq N$,
    dann ist die Reihe konvergent.
    \item[\rlap{\raisebox{0.3ex}{\hspace{0.4ex}\tiny \ding{52}}}$\square$]
    Falls die Reihe konvergiert, so folgt $\forall \epsilon > 0, \exists N \geq 1$ gilt 
    $\sum_{k = n}^{n+100} |a_k| < \epsilon, \forall n \geq N$.
    \item[$\square$]
    Falls die Reihe $\Reihe \sin(a_k)$ absolut konvergiert, so konvergiert die ursprüngliche Reihe.
\end{itemize}

$\sum_{n = 1}^{\infty} c_n$ und $\alpha > 0$, $a_n = c_n\alpha^n$
und $b_n = nc_n\alpha^{n-1}$.
\begin{itemize}
    \item[$\square$]
    $\limsup_{n \to \infty} |a_n|^{1/n} > \limsup_{n \to \infty} |b_n|^{1/n}$
    \item[$\square$]
    $\limsup_{n \to \infty} |a_n|^{1/n} < \limsup_{n \to \infty} |b_n|^{1/n}$
    \item[\rlap{\raisebox{0.3ex}{\hspace{0.4ex}\tiny \ding{52}}}$\square$]
    $\limsup_{n \to \infty} |a_n|^{1/n} = \limsup_{n \to \infty} |b_n|^{1/n}$
\end{itemize}

Die Potenzreihe $\sum_{n = 1}^\infty \frac{2^n}{n\sqrt{n}} \cdot z^n$ ist auf dem Rand ihres 
Konvergenzkreises...
\begin{itemize}
    \item[\rlap{\raisebox{0.3ex}{\hspace{0.4ex}\tiny \ding{52}}}$\square$]
    überall absolut konvergent
    \item[$\square$]
    überall konvergent, aber nicht absolut konvergent
    \item[$\square$]
    überall konvergent, aussert in unendlich vielen Punkten
    \item[$\square$]
    nirgendwo konvergent
\end{itemize}

Sei $\phi$ eine Abbildung einer Reihe $\Reihe a_k$ und $b_n = a_{\phi(n)}$. Dann gilt:
\begin{itemize}
    \item[$\square$]
    Wenn die Reihe konvergiert und $\phi$ surjektiv ist, dann ist die Reihe mit $b_n$ auch konvergent.
    \item[$\square$]
    Wenn die Reihe konvergiert und $\phi$ injektiv ist, dann ist die Reihe mit $b_n$ auch konvergent.
    \item[$\square$]
    Wenn die Reihe absolut konvergiert und $\phi$ surjektiv ist, dann ist die Reihe mit $b_n$ auch konvergent.
    \item[\rlap{\raisebox{0.3ex}{\hspace{0.4ex}\tiny \ding{52}}}$\square$]
    Wenn die Reihe absolut konvergiert und $\phi$ injektiv ist, dann ist die Reihe mit $b_n$ auch konvergent.
\end{itemize}

Sei $\Reihe a_k$ absolut konvergent, dann gilt für die Reihe $\Reihe |a_k|^2$
\begin{itemize}
    \item[$\square$]
    konvergiert nicht unbedingt
    \item[\rlap{\raisebox{0.3ex}{\hspace{0.4ex}\tiny \ding{52}}}$\square$]
    konvergiert immer absolut
    \item[$\square$]
    konvergiert immer, aber nicht absolut
\end{itemize}

\Title{Funktionen}

Seien $f : \R \to \R$ und $h : \R \to [0, \infty[$, so dass $\lim_{x \to \infty}f(x) = - -\infty, \lim_{x \to \infty} h(x) = 0$.
\begin{itemize}
    \item[\rlap{\raisebox{0.3ex}{\hspace{0.4ex}\tiny \ding{52}}}$\square$]
    Ist $\lim_{x \to \infty} f(x)h(x) = 0$ möglich?
    \item[$\square$]
    Ist $\lim_{x \to \infty} f(x)h(x) = \infty$ möglich?
    \item[\rlap{\raisebox{0.3ex}{\hspace{0.4ex}\tiny \ding{52}}}$\square$]
    Ist $\lim_{x \to \infty} f(x)h(x) = - \infty$ möglich?
    \item[\rlap{\raisebox{0.3ex}{\hspace{0.4ex}\tiny \ding{52}}}$\square$]
    Ist $\lim_{x \to -\infty} f(x)h(x) = \infty$ möglich?
\end{itemize}

Seien $f,g$ monoton wachsende Funktionen:
\begin{itemize}
    \item[$\square$]
    $f \cdot g$ ist monoton wachsend. 
    \item[$\square$]
    $\frac{f}{g}$ ist monoton wachsend für $g(x) \neq 0$.
    \item[$\square$]
    $\frac{f}{g}$ oder $\frac{g}{f}$ ist monoton wachsend für $f(x), g(x) \neq 0$.
\end{itemize}

Seien $f : X \to Y, g : Y \to Z$ Funktionen, so dass $g \circ f : X \to Z$ eine Bijektion ist. Welche
Aussage ist richtig?
\begin{itemize}
    \item[$\square$]
    $f$ ist injektiv, $g$ ist injektiv
    \item[\rlap{\raisebox{0.3ex}{\hspace{0.4ex}\tiny \ding{52}}}$\square$]
    $f$ ist injektiv, $g$ ist surjektiv
    \item[$\square$]
    $f$ ist surjektiv, $g$ ist injektiv
    \item[$\square$]
    $f$ ist surjektiv, $g$ ist surjektiv
\end{itemize}

Sei $f : [a,b] \to \R$ eine Funktion. Sei $f_n$ eine Funktionenfolge, die gleichmässig gegen $f$
konvergiert. Es gilt:
\begin{itemize}
    \item[$\square$]
    Sei $x_0 \in ]a,b[$. Falls $f_n$ für alle $n \geq 1$ in $x_0$ differnzierbar ist, so ist $f$ in 
    $x_0$ differnzierbar.
    \item[\rlap{\raisebox{0.3ex}{\hspace{0.4ex}\tiny \ding{52}}}$\square$]
    Falls $f_n$ für alle $n \geq 1$ in $x_0$ auf $[a,b]$ beschränkt ist, so existiert für jede Partition
    $P$ der Grenzwert der Untersumme $\Limit s(f_n, P) = s(f,P)$.
    \item[\rlap{\raisebox{0.3ex}{\hspace{0.4ex}\tiny \ding{52}}}$\square$]
    Falls $f_n$ für alle $n \geq 1$ in $x_0$ stetig ist, so ist $f$ gleichmässig stetig. 
    \item[\rlap{\raisebox{0.3ex}{\hspace{0.4ex}\tiny \ding{52}}}$\square$]
    Falls $f_n$ für alle $n \geq 1$ in $x_0$ konvex ist, so ist $f$ konvex.
\end{itemize}

Sei $D \subseteq \R$. Dann gilt:
\begin{itemize}
    \item[$\square$]
    Falls $D = \emptyset$, dann besitzt $D$ einen Häufungspunkt.
    \item[\rlap{\raisebox{0.3ex}{\hspace{0.4ex}\tiny \ding{52}}}$\square$]
    Falls $D \subseteq E$ und $x_0$ ein Häufungspunkt von $D$, so ist $x_0$ auch ein Häufungspunkt von $E$.
    \item[$\square$]
    Falls $F \subseteq D$ und $x_0$ ein Häufungspunkt von $D$, so ist $x_0$ auch ein Häufungspunkt von $F$.
    \item[\rlap{\raisebox{0.3ex}{\hspace{0.4ex}\tiny \ding{52}}}$\square$]
    Falls $D$ endlich ist, gibt es keinen Häufungspunkt.
    \item[$\square$]
    Falls $D$ unendlich ist, gibt es mindestens einen Häufungspunkt.
\end{itemize}

Sei $f_n : D \to \R$. Dann gilt:
\begin{itemize}
    \item[\rlap{\raisebox{0.3ex}{\hspace{0.4ex}\tiny \ding{52}}}$\square$]
    $f_n$ konvergiert punktweisen, wenn $|f(x) - f_n(x)| \to 0,$ für $n \to \infty$.
    \item[$\square$]
    $f_n$ konvergiert gleichmässig, wenn $\forall \epsilon > 0, \forall x \in D, \exists N > 0$, so dass $|f(x) - f_n(x)| < \epsilon$.
    \item[\rlap{\raisebox{0.3ex}{\hspace{0.4ex}\tiny \ding{52}}}$\square$]
    Gleichmässige konvergenz von $f_n$ impliziert punktweise konvergenz von $f_n$.
\end{itemize}

Welche Aussagen sind richtig?
\begin{itemize}
    \item[$\square$]
    $f : [0,1] \to \R$ beschränkt $\Rightarrow$ $f$ monoton.
    \item[$\square$]
    $f : [0,1] \to \R$ strikt monoton wachsend $\Rightarrow$ $f$ stetig.
    \item[$\square$]
    $f : ]0,1] \to \R$ monoton $\Rightarrow$ $f$ beschränkt.
    \item[\rlap{\raisebox{0.3ex}{\hspace{0.4ex}\tiny \ding{52}}}$\square$]
    $f : [0,1] \to \R$ monoton $\Rightarrow$ $f$ beschränkt.
\end{itemize}

Sei $f : [0,1] \to [0,1]$ stetig und nicht konstant. Dann gilt:
\begin{itemize}
    \item[\rlap{\raisebox{0.3ex}{\hspace{0.4ex}\tiny \ding{52}}}$\square$]
    Es gibt einen Punkt $x \in [0,1]$, so dass $f(x) = x$.
    \item[$\square$]
    Es gibt einen Punkt $x \in [0,1]$, so dass $f(x) = 1$.
    \item[$\square$]
    $f$ hat eine eindeutige Maximalstelle.
    \item[\rlap{\raisebox{0.3ex}{\hspace{0.4ex}\tiny \ding{52}}}$\square$]
    Das Bild $f([0,1]) \subset [0,1]$ ist ein abgeschlossenes Intervall. D.h. es gibt $a,b \in [0,1]$
    mit $a < b$, so dass $f([0,1]) = [a,b]$.
\end{itemize}

Sei $f : \R \to \R$ stetig bei $x_0 = 0$, mit $f(x_0) > 0$. Dann gilt:
\begin{itemize}
    \item[$\square$]
    Es gilt $f(x) > 0$ für alle $x \in \R$.
    \item[\rlap{\raisebox{0.3ex}{\hspace{0.4ex}\tiny \ding{52}}}$\square$]
    Es existiert $\epsilon, \delta > 0$, so dass $f(x) > \epsilon$ für alle $x \in (- \delta, \delta)$ gilt.
\end{itemize}

Sei $f : D \to \R$, welche Aussage ist äquivalent zur Stetigkeit von $f$.
\begin{itemize}
    \item[\rlap{\raisebox{0.3ex}{\hspace{0.4ex}\tiny \ding{52}}}$\square$]
    $\forall x \in D, \forall \epsilon > 0, \exists \delta > 0$, so dass für alle $z \in D$ gilt: $z \in (x - \delta, x + \delta) \Rightarrow f(z) \in (f(x) - \epsilon, f(x) + \epsilon)$.
    \item[$\square$]
    $\forall x \in D, \exists \delta > 0, \forall \epsilon > 0$, so dass für alle $z \in D$ gilt: $|z - x| < \delta \Rightarrow |f(z) - f(x)| < \epsilon$.
    \item[$\square$]
    $\forall \epsilon > 0, \exists \delta > 0$, so dass für alle $x,z \in D$ gilt: $|z - x| < \delta \Rightarrow |f(z) - f(x)| < \epsilon$.
\end{itemize}

Sei $f : \R \to \R$ stetig. Dann gilt:
\begin{itemize}
    \item[$\square$]
    $f$ hat eine Maximalstellt
    \item[\rlap{\raisebox{0.3ex}{\hspace{0.4ex}\tiny \ding{52}}}$\square$]
    Wenn $f(1) = 2$ und $f(3) = -1$, dann gibt es ein $t \in ]1,3[$ mit $f(t) = 0$.
    \item[\rlap{\raisebox{0.3ex}{\hspace{0.4ex}\tiny \ding{52}}}$\square$]
    Wenn $\lim_{t \to \infty} f(t) = 0$ gilt, dann ist $f$ auf $[0, \infty[$ begrenzt.
    \item[$\square$]
    Wenn $f$ begrenzt ist, dann existiert $\lim_{t \to \infty} f(t)$.
\end{itemize}

Seien $(a_n), (b_n), (c_n)$ Folgen mit $c_n = a_n + b_n$. Dann gilt:
\begin{itemize}
    \item[$\square$]
    Falls $\Limit c_n$ existiert, existieren $\Limit a_n$ und $\Limit b_n$ und es gilt
    $\Limit c_n = \Limit a_n + \Limit b_n$.
    \item[\rlap{\raisebox{0.3ex}{\hspace{0.4ex}\tiny \ding{52}}}$\square$]
    Falls $\Limit c_n$ und $\Limit b_n$ existieren, existiert $\Limit a_n$  und es gilt
    $\Limit a_n = \Limit c_n - \Limit b_n$.
\end{itemize}

Seien $a < b$, $g : \R \to \R$ beschränkt und $f : [a, b] \to \R$ brschränkt mit $f(a) < f(b)$. Dann gilt:
\begin{itemize}
    \item[$\square$]
    Falls für jedes $c \in [f(a), f(b)]$ ein $x \in [a,b]$ mit $f(x) = c$ existiert, so folgt, dass 
    $f$ stetig ist.
    \item[$\square$]
    Falls $g \circ f$ und $g$ stetig sind, so ist auch $f$ stetig.
    \item[\rlap{\raisebox{0.3ex}{\hspace{0.4ex}\tiny \ding{52}}}$\square$]
    Falls $f$ stetig ist, gibt es $x_0 \in [a,b]$, so dass $\int_a^b x f(x)dx = \frac{f(x_0)}{2}(b^2 - a^2)$.
\end{itemize}

Sei $f : [0, \infty[ \to \R$ stetig. Dann gilt:
\begin{itemize}
    \item[\rlap{\raisebox{0.3ex}{\hspace{0.4ex}\tiny \ding{52}}}$\square$]
    Es gibt eine differnzierbare Funktion $F : ]0, \infty[ \to \R$ mit $F'(t) = f(t)$ für alle $t \in ]0 \infty[$.
    \item[\rlap{\raisebox{0.3ex}{\hspace{0.4ex}\tiny \ding{52}}}$\square$]
    Wenn $f(t) \geq 0$ für alle $t \in [0,1]$ und $\int_0^1 f(t)dt = 0$, dann gilt $f(t) = 0$ für alle
    $t \in [0,1]$.
    \item[$\square$]
    Wenn $|f(t)| \leq \frac{1}{1+ t}$ für alle $t \in [0, \infty[$, dann existiert $\int_0^\infty f(t)dt$.
    \item[$\square$]
    Wenn $\sum_{n = 1}^\infty f(x)$ konvergiert, dann existiert $\int_0^\infty f(t)dt$.
\end{itemize}

\Title{Ableiten und Integrieren}

Seien $f(x) = \sin(x) \cdot e^{-1/x^2}$ mit $f(0) = 0$ und $a_k = \int_{2k\pi}^{2(k+1)\pi} f(x)dx$. Dann gilt:
\begin{itemize}
    \item[$\square$]
    $f$ ist stetig aber nicht glatt.
    \item[\rlap{\raisebox{0.3ex}{\hspace{0.4ex}\tiny \ding{52}}}$\square$]
    $f$ besitzt unendlich viele lokale Minimalstellen.
    \item[$\square$]
    $f$ besitzt ein lokales Maximum in $x = 0$.
    \item[$\square$]
    $a_1 > 0$ 
    \item[\rlap{\raisebox{0.3ex}{\hspace{0.4ex}\tiny \ding{52}}}$\square$]
    $\lim_{k \to \infty} a_k = 0$.
\end{itemize}

Sei $f$ differenzierbar. Dann gilt:
\begin{itemize}
    \item[\rlap{\raisebox{0.3ex}{\hspace{0.4ex}\tiny \ding{52}}}$\square$]
    $f$ ist stetig.
    \item[$\square$]
    $f'$ ist stetig.
    \item[$\square$]
    Falls $f'(x) = 0, \forall x \in D$, dann existiert $c \in \R, f(x) = c, \forall x \in D$.
    \item[\rlap{\raisebox{0.3ex}{\hspace{0.4ex}\tiny \ding{52}}}$\square$]
    Falls $D = (a,b), a < b$, dann ist obige Aussage wahr.
\end{itemize}

Sei $f$ eine ungerade Funktion. Dann gilt:
\begin{itemize}
    \item[$\square$]
    $f^{(i)}(0) = 0$ für $i$ ungerade
    \item[$\square$]
    $f^{(i)}(0) \neq 0$ für $i$ ungerade
    \item[\rlap{\raisebox{0.3ex}{\hspace{0.4ex}\tiny \ding{52}}}$\square$]
    $f^{(i)}(0) = 0$ für $i$ gerade
    \item[$\square$]
    $f^{(i)}(0) \neq 0$ für $i$ gerade
\end{itemize}

Welche der Implikationen sind wahr?
\begin{itemize}
    \item[\rlap{\raisebox{0.3ex}{\hspace{0.4ex}\tiny \ding{52}}}$\square$]
    $f$ differenzierbar $\Rightarrow$ $f$ stetig $\Rightarrow$ $f$ integrierbar
    \item[$\square$]
    $f$ integrierbar $\Rightarrow$ $f$ differenzierbar $\Rightarrow$ $f$ stetig
    \item[$\square$]
    $f$ stetig $\Rightarrow$ $f$ differenzierbar $\Rightarrow$ $f$ integrierbar
    \item[$\square$]
    $f$ integrierbar $\Rightarrow$ $f$ stetig $\Rightarrow$ $f$ differenzierbar
\end{itemize}

Sei $h(x) = g(f(x))$ und $f,g$ sind zwei nicht-negative Funktionen. 
\begin{itemize}
    \item[$\square$]
    Wenn $f$ und $g$ streng konvex sind, ist auch $h$ konvex.
    \item[\rlap{\raisebox{0.3ex}{\hspace{0.4ex}\tiny \ding{52}}}$\square$]
    Wenn $\int_0^\infty f(x)dx, \int_0^\infty g(x)dx$ konvergieren, so konvergiert auch $\int_0^\infty h(x)dx$.
\end{itemize}

Welche Aussage stimmt?
\begin{itemize}
    \item[$\square$]
    Seien $f_n : [a,b] \to \R$ beschränkt und integrierbar. Dann gilt für $f(x) := \Limit f_n(x) : \sum_{n = 0}^\infty \int_a^bf_n(x)dx = \int_a^b \sum_{n = 0}^\infty f_n(x)dx$.
    \item[\rlap{\raisebox{0.3ex}{\hspace{0.4ex}\tiny \ding{52}}}$\square$]
    Seien $f_n$ wie oben, aber zusätzlich konvergieren sie gleichmässig. Dann gilt: $\Limit \int_a^b f_n(x)dx = \int_a^bf(x)dx$.
    \item[$\square$]
    Die Umkehrung der zweiten Aussage ist wahr.
\end{itemize}


\Title{Weiter Aufgaben}

\Title{Folgen und Reihen}

\textbf{Ex. Sei $(a_n)$ eine induktive Folge mit $a_1 = \sqrt{2}, a_{n+1} = \sqrt{2 + a_n}$. Beweise, dass die Folge durch $2$ von oben beschränkt ist und berechne den Grenzwert.} 

Wir wissen, dass die Wurzelfunktion monoton wachsend ist. Nun machen wir einen Induktionsbeweis. Verankerung:
für $n = 1, a_1 = \sqrt{2} \leq 2$. Induktionsschritt: $a_{n+1} = \sqrt{2 + a_n} \leq \sqrt{2 + 2} = 2$.
Somit haben wir die Beschränktheit gezeigt, nun beweisen wir noch, dass die Folge monoton wachsend ist.
Verankerung: für $n = 2, a_2 = \sqrt{2 + \sqrt{2}} \geq \sqrt{2} = a_1$. Induktionsschritt:
$a_{n+1} = \sqrt{2 + a_n} \geq \sqrt{2 + a_{n-1}} = a_n$. Da die Folge beschränkt und monoton wachsend ist,
folgern wir, dass ein Grenzwert existiert. $a = \Limit a_n = \sqrt{2 + a} \Rightarrow a^2 - 2 - a = 0$,
da $a_n > 0$ kommt nur die Lösung $a = 2$ in frage. \medskip

\textbf{Ex. Zeige, dass $\Limit \prod_{k = 1}^n(1 + \frac{t}{k^2})$ für jedes $t$ existiert.} 

Wir nehmen $N \in \N, \frac{|t|}{N^2} < 1$. Da $\prod_{k = 1}^n(1 + \frac{t}{k^2}) = \prod_{k = 1}^N(1 
+ \frac{t}{k^2}) \prod_{k = N + 1}^n(1 + \frac{t}{k^2})$ gilt, reich es wenn wir zeigen, dass 
$\Limit \prod_{k = N+1}^n(1 + \frac{t}{k^2}) = \Limit \prod_{k = 1}^n(1 + \frac{t}{(k + N)^2}) = a_n$ existiert.
Falls $t \geq 0$ ist die Folge $(a_n)$ steigend, sonst ist sie fallend. Weiter ist sie monoton, deswegen
existiert ein Grenzwert, wenn sie beschränkt ist. Es folgt $0 \leq a_n \leq \prod_{k = 1}^n(e^{\frac{t}{k^2}})
\leq e^{|t|(\sum_{k = 1}^\infty\frac{1}{k^2})} < \infty$. Weshalb der Grenzwert existiert. \medskip

\textbf{Ex. Zeige mit der Def. des Grenzwertes, dass $\Limit (\sqrt{2n + \cos(n)} - \sqrt{2n + 1}) = 0$.}

Wir formen zuerst um und schätzen ab: $|\sqrt{2n + \cos(n)} - \sqrt{2n + 1}| = \frac{|\cos(n)-1|}{|\sqrt{2n + \cos(n)} + \sqrt{2n + 1}|}
\leq \frac{2}{\sqrt{n + 1}}$. Nun verwenden wir die definition des Grenzwertes: $|\frac{2}{\sqrt{n + 1}} - 0| < \epsilon 
\Rightarrow n > \frac{4 - \epsilon^2}{2 \epsilon^2}$. Somit haben wir gezeigt, dass für ein beliebiges
$\epsilon$ ein $n$ existiert, so dass $|a_k - 0| < \epsilon, \forall k > n$. \medskip

\textbf{Ex. Berechne den Grenzwert $\lim_{x \to 1^+} \frac{\log(x)-\sin(x\pi)}{\sqrt{x-1}}$.}

Wenn $|x-1|$ klein genug ist können wir die Taylor Reihe verwenden. Dann gilt: $\frac{\log(x)-\sin(x\pi)}{\sqrt{x-1}} =
\frac{(1 + \pi)(x-1) + \mathcal{O}(|x-1|^2)}{\sqrt{x-1}} = (1 + \pi)\sqrt{x-1} + \mathcal{O}(|x-1|^\frac{3}{2})$.
Somit ist $\lim_{x \to 1^+} \frac{\log(x)-\sin(x\pi)}{\sqrt{x-1}} = 0$. \medskip

\textbf{Ex. Zeige, dass für alle $t \in \R$ $\Limit(\cos(\frac{t}{\sqrt{n}}))^n = e^{-\frac{t^2}{2}}$.} 

Wir verwenden wieder die Taylorreihen für $\cos$ und $\log$. $\cos(\frac{t}{\sqrt{n}})^n = \exp(n \log(\cos(\frac{t}{\sqrt{n}}))) = 
\exp(n \log(1 - \frac{t^2}{2n} + \mathcal{O}(\frac{1}{n^2}))) = \exp(n (- \frac{t^2}{2n} + \mathcal{O}(\frac{1}{n^2}))) = 
\exp(- \frac{t^2}{2} + \mathcal{O}(\frac{1}{n})) \to^{n \to \infty} \exp(-\frac{t^2}{2})$. \medskip

\textbf{Ex. Berechne $\lim_{x \to 0} \frac{\sin(x)}{x}$.} 

Wir verwenden die Potenzreihe Darstellung des Sinus
$\frac{1}{x} \Reihe (-1)^k \frac{x^{2k+1}}{(2k+1)!} = \Reihe (-1)^k \frac{x^{2k}}{(2k+1)!}$.
Zudem zeigen wir mit dem Quotientenkriterium,
dass diese Potenzreihe für jedes $x \in \R$ absolut konvergent ist. Da der Wert der Potenzreihe für 
$x = 0$ gegen 1 konvergiert, schliessen wir $\lim_{x \to 0} \frac{\sin(x)}{x} = 1$.
\medskip

\textbf{Ex. Für welche $x$ konvergiert die Reihe $\Reihe |x|^{n!}$?} 

Falls $|x| \geq 1$ sehen wir, dass $\Limit |x|^{n!} \neq 0 $, also muss die Reihe divergieren.
Falls $x < 1$ ist $\Reihe |x|^{n!} \leq \Reihe |x|^n$ und da dies eine Geometrische Reihe mit Basis < 1
ist, konvergiert die Reihe. \medskip


\textbf{Ex. Stelle $(1+e^x)^3$ als Potenzreihe dar.} 

Es gilt $\Reihe b_k = b, \Reihe c_k = c$ und $\Reihe b_k + c_k = b + c$. Also können wir die
Potenzreihe in einzelne Summanden aufteilen: $(1+e^x)^3 = 1 + 3e^x + 3e^{2x} + e^{3x}
= 1 + \Sum_{n = 0}^\infty 3 \frac{x^n}{n!} + 3 \frac{2^nx^n}{n!} + \frac{3^nx^n}{n!} = 
1 + \Sum_{n = 0}^\infty \frac{3 + 3 \cdot ^n + 3^n}{n!}x^n$. Nun müssen wir nur noch den Fall $n = 0$
und $n > 0$ unterscheiden. \medskip

\textbf{Ex. Zeige, dass das Cauchy Produkt der beiden divergenten Reihen $2 + 2 + 2^2 + ...$ und $-1 + 1 + 1 + ...$ absolut konvergiert.} 

Wir berechnen $\sum_{n = 0}^\infty(\sum_{j = 0}^n a_{n - j}b_j) = a_0 b_0 + \sum_{n = 1}^\infty(\sum_{j = 0}^n a_{n - j}b_j)
= \sum_{n = 0}^\infty(-2^n + 2^{n-1} + ... + 2^1 + 2) = -2 + \sum_{n = 1}^\infty(1 - (1 - 2^n) - 2^n) = \Reihe c_k$
mit $c_0 = -2$ und $c_n = 0$. Offensichtlich konvergiert $c_n$ absolut. \medskip

\textbf{Ex. Sei $\Reihe a_k$ absolut konvergent und $\Reihe b_k$ nur konvergent. Folgt daraus, dass $\Reihe b_k \sin(a_k)$ konvergiert?} 

Zuerst zeigen wir, dass $\Limit b_n = 0$. Da $\Reihe b_k$ konvergiert, wissen wir dass $\Reihe c_k$
konvergiert, wobei $c_1 = 0$, $c_k = b_{k-1}$ und $\Limit \Sum_{k = 1}^n b_k = \Limit \Sum_{k = 1}^n c_k$
$\to$ $\Limit b_k = \Limit (\Sum_{k = 1}^n b_k - \Sum_{k = 1}^n c_k) = 0$. Nun verwenden wir die
Abschätzung $|\sin(x)| \leq |x|$, woraus folgt, dass $\Reihe \sin(a_k)$ absolut konvergiert. Wir wissen
dass die Folge $(b_k)$ einen oberen Grenzwert $C > 0$ hat, deshalb gilt: $|b_n \sin(a_n)| \leq C|a_n|$,
was den Beweis abschliesst. \medskip

\Title{Funktionen}

\textbf{Ex. Zeige für alle $k \in \N$ und $0 \leq x \leq \sqrt{(4k+5)(4k+6)}$ gilt: $\cos(x) \geq 1 - \frac{x^2}{2!} + \frac{x^4}{4!} - \dots - \frac{x^{2(2k+1)}}{2(2k+1)!}$.} 

Der Kosinus konvergiert gleichmässig und $\cos(x) - (1 - \frac{x^2}{2!} + \frac{x^4}{4!} - \dots - \frac{x^{2(2l+1)}}{2(2l+1)!}) 
= \Reihe \frac{x^{4(k + l)}}{(4k + 4l)!} - \frac{x^{4(k + l) + 2}}{(4k + 4l + 2)!}$. Da 
$x \leq \sqrt{(4(k+l) + 1)(4(k+l) + 2)} \Rightarrow \frac{x^{4(k + l)}}{(4k + 4l)!} - \frac{x^{4(k + l) + 2}}{(4k + 4l + 2)!} \geq 0$,
können wir schliessen, dass $0 \leq x \leq \sqrt{(4k+5)(4k+6)} \Rightarrow \frac{x^{4(k + l)}}{(4k + 4l)!} - \frac{x^{4(k + l) + 2}}{(4k + 4l + 2)!} \geq 0, \forall l \geq 1$.
\medskip

\textbf{Ex. Konvergiert $f_n(x) = x \to 2nx$ falls $0 \leq x < 1/2n$, $2-2nx$ falls $1/2n \leq x < 1/n$, $0$ falls $1/n \leq x \leq 1$ gleichmässig?} 

Nein, $f_n(x)$ konvergiert punktweise gegen $f(x) = 0$. Für $x= 0$, gilt $f_n(x) = 0$ für jedes $n$. Für
$x > 0$ wählen wir $n_0 \in \N$ so dass $x > 1/n_0$. Dann gilt $x > 1/n$ für alle $n \geq n_0$. Also ist 
$f_n(x) = 0$ für $n \geq n_0$. Wir behaupten nun, dass $f_n(x)$ nicht gleichmässig konvergiert.
Es gilt für jedes $n$, dass $f_n(\frac{1}{2n}) = 1$ und damit $f_n(\frac{1}{2n}) - f(\frac{1}{2n}) = 1$.
Damit gilt aber auch $\sup_{x\in[0,1]} |f_n(x) - f(x)| \geq |f_n(\frac{1}{2n}) - f(\frac{1}{2n})| = 1 \neq 0 $
und die Funktionenfolge konvergiert nicht gleichmässig. \medskip

\textbf{Ex. Sei $x_0$ ein Häufungspunkt von $D \cap ]x_0 , \infty[, D \cap ]- \infty , x_0[$ mit der Eigenschaft
$\lim_{x \to x_0^+} f(x) = +\infty$, $\lim_{x \to x_0^-} f(x) = -\infty$. Zeige, dass $\lim_{x \to x_0} \frac{1}{1 + f(x)^2} = 0$.} 

Wir verwenden die Definition des Häufungspunktes. Es gilt $\frac{1}{1 + f(x)^2} \leq \frac{1}{1 + \frac{1}{\epsilon^2}} = \frac{\epsilon^2}{\epsilon^2 + 1}$.
Wir wählen nun $\delta'' = \text{min}(\delta, \delta')$. Dann erhalten wir $|\frac{1}{1 + f(x)^2} - 0| \leq \frac{\epsilon^2}{\epsilon^2 + 1} = \epsilon' $.
Nach Definition haben wir somit gezeigt, dass der Grenzwert 0 ist. \medskip

\textbf{Ex. Bestimme die Konstanten $a, b$ so, dass die Funktione $f : x \to x^2-ax+b$ falls $x \leq -1$, $(a + b)x$ falls $-1 < x < 1$, $x^2 + ax - b$ falls $x \geq 1$ auf ganz $\R$ stetig ist.}
Die einzelnen Funktionen sind stetig, daher müssen wir nur noch die Punkte $\pm 1$ überprüfen. Es muss gelten
$\lim_{x \to -1^+} (a + b)x = 1 + a + b$ und $\lim_{x \to 1^-} (a + b)x = 1 + a - b$. Wenn wir diese 
Gleichungen lösen, erhalten wir $a = -1$ und $b = \frac{1}{2}$. \medskip

\textbf{Ex. Sei $f : x \to \frac{x^5 - 7x^2 + 1}{x^8 + 1}$, zeige, dass $f$ auf ganz $\R$ stetig ist und min. eine Nullstelle besitzt.} 

Arithmetische Operationen sind stetig auf $\R$ und die Verkettung von stetigen Funktionen ist ebenfalls stetig.
Somit wissen wir $f$ ist stetig auf ganz $\R$. Es gilt $f(-1) < 0$ und $f(2) > 0$, mit dem Zwischenwertsatz
folgt nun dass $f$ eine Nullstelle in $[-1, 2]$ hat. \medskip

\textbf{Ex. Sei $f(x) = x$ für $x \in \Q$ und $f(x) = 1 - x$ sonst. Zeige, dass $x_0 = \frac{1}{2}$ der einzige Stetigkeitspunkt ist.} 

Wir nehemn an $x_0 \neq 1/2$ und machen eine Fallunterscheidung. Zuerst nehmen wir an $x_0 \in \R \backslash \Q$
ist ein Stetigkeitspunkt. Wir definieren $(y_n) = \frac{\lfloor n x_0\rfloor}{n}$, somit ist $\Limit y_n = x_0$ 
und $\Limit f(y_n) = x_0 \neq 1-x_0 = f(x_0)$. Somit ist $f$ für alle $x_0 \in \R \backslash \Q$ nicht stetig. 
Nun nehmen wir an $x_0 \in \Q \backslash \{1/2\}$. Wir definieren $(y_n) = x_0 + \frac{\sqrt{2}}{n} \in \R \backslash \Q$ 
$\Limit y_n = x_0$ und $\Limit f(y_n) = 1-x_0 \neq x_0 = f(x_0)$. Somit ist $f$ für alle 
$x_0 \in \Q \backslash \{1/2\}$ nicht stetig. Nun müssen wir noch zeigen, dass $f$ in $x_0$ stetig ist.
D.h. $\forall \epsilon \exists \delta |x - x_0| < \delta \Rightarrow |f(x) - f(x_0)| < \epsilon$.
Merke $x_0 \in \Q$ nun wählen wir $\delta = \epsilon$ und erhalten $|f(x) - f(x_0)| = |f(x) - 1/2| = |x - 1/2|
= |x - x_0| < \delta = \epsilon$. Daher ist $f$ stetig in $x_0$. \medskip

\textbf{Ex. Sei $f : [0, \ln(2)] \to \R$ eine stetige Funktion. Zeigen Sie, dass es $\eta \in [0, \ln(2)],$ $f(\eta) = \frac{1}{e^2 - e}\int_0^{\ln(2)}e^{e^x}e^xf(x)dx$ gibt.}

Nach dem Min-Max Satz hat $f$ ein Maximum $b$ und ein Minimum $a$, so dass $f[0, \ln(2)] \to [f(a), f(b)]$. Daher gilt: $f(a) \int_0^{\ln(2)}e^{e^x}e^xdx \leq \int_0^{\ln(2)}e^{e^x}e^xf(x)dx \leq f(b)\int_0^{\ln(2)}e^{e^x}e^xdx$. Nun rechnen wir $\int_0^{\ln(2)}e^{e^x}e^xdx = e^2-e$ aus. Wenn wir den Zwischenwertsatz anwenden, erhalten wir somit es existiert ein $\eta \in [0, \ln(2)]$, so dass $f(a) \leq f(\eta) = \frac{1}{e^2 - e} \int_0^{\ln(2)}e^{e^x}e^xf(x)dx \leq f(b)$.
\medskip

\textbf{Ex. Sei $f_n : [0,1] \to \R, x \to \frac{nx}{(n^2x^2+1)^2}$. Zeige, dass $f(x) = \Limit f_n(x)$. Konvergiert die Funktion gleichmässig? Gilt $\Limit \int_0^1 f_n(x)dx = \int_0^1f(x)dx$?}

Es gilt $\Limit f_n(x) = 0 = f(x)$. Die Konvergenz ist nicht gleichmässig, da $\sup_{0 \leq x \leq 1}|f_n(x)| \geq \frac{1}{4}, \forall n \in \N$. Weiter gilt $\Limit \int_0^1f_n(x)dx = \int_0^1f(x)dx = 0$, denn $\int_0^1f_n(x)dx = -\frac{1}{2n} \frac{1}{n^2x^2 + 1}|_0^1 \to 0$.
\medskip

\Title{Ableiten und Integrieren}

\textbf{Ex. Zeige, dass $f : x \to x + e^x$ bijektiv von $\R$ nach $\R$ ist und die Inverse auf ganz
$\R$ differenzierbar ist.}

$f'(x) = 1 + e^x > 0, \forall x \in \R$, somit ist $f$ streng monoton wachsend und umkehrbar. Dazu gilt
$\lim_{x \to - \infty}f(x) = - \infty$ und $\lim_{x \to \infty} f(x)= \infty$ somit ist $f$ bijektiv von 
$\R$ zu $\R$. \medskip

\textbf{Ex. Zeige, dass $1 + x \leq e^x, \forall x \in \R$.}

Wir definieren $f : x \to e^x - 1 - x$. Da die Exponentialfunktion schneller als jedes Polynom wächst,
gilt $\lim_{x \to - \infty} f(x) = \lim_{x \to \infty} f(x) = \infty$. Da $f$ stetig ist, nimmt sie
ein Minimum an. Da $f'(x) = 0$ genau für $x = 0$, ist dies eine globale Minimalstellen. Somit folgt 
$0 \leq f(x) \Rightarrow 1 + x \leq e^x$. \medskip

\textbf{Ex. Berechne die Minimal-/Maximalstellen von $f : x \to \sqrt{1 + x} - \frac{1}{2}\sqrt{x}, x \in [0,5]$.} 

$f'(x) = \frac{1}{4\sqrt{x(x+1)}}(2\sqrt{x}-\sqrt{1+x})$, also ist $f' > 0$ genau dann wenn, $x > \frac{1}{3}$.
Dann ist $f$ monoton fallend auf $[0,\frac{1}{3}]$ und stark monoton wachsend auf $[\frac{1}{3},5]$.
Wenn wir nun noch $f(5) > f(0)$ vergleichen, sehen wir das die Maximalstelle bei $5$ und die Minimalstelle 
bei $\frac{1}{3}$ liegt. \medskip

\textbf{Ex. Differenziere $f(x) = \int_{\cos(x)}^{e^x}\cos(t)dt$.}

$f(x) = F(e^x) - F(\cos(x)) \Rightarrow f'(x) = F'(e^x) - F'(\cos(x))$ \\
Wir wissen, die Ableitung eines Integrals ist die innere Funktion, dazu kommt nun noch innere Ableitung
und wir erhalten:
$f'(x) = \cos(e^x) \cdot e^x + \cos(\cos(x)) \cdot \sin(x)$ \medskip

\textbf{Ex. Sei $f$ differenzierbare mit $f(x_0) \neq 0$ für mindestens ein $x_0 \in \R$. Weiter gilt
$f(x + y) = f(x)f(y)$. Zeige $f(0) = 1$.}

Sei $b := f(x_0)$, so dass $b \neq 0$. Dann gilt $bf(0) = f(x_0)f(0) = f(x_0 + 0) = f(x_0) = b$,
also ist $f(0) = \frac{bf(0)}{b} = \frac{b}{b} = 1$. \medskip

\textbf{Ex. Sei $f(x) = x\sin(\frac{1}{x})$ mit $f(0) = 0$. Zeige, dass $f$ in $0$ nicht differenzierbar ist.} 

Es reicht eine Folge $(y_n)$ zu finden, die gegen $0$ strebt sodass $\frac{f(y_n)-f(0)}{y_n}$ nicht
konvergiert. Dafür wählen wir $y_n = \frac{2}{\pi n}$. Es gilt $\frac{f(y_n)-f(0)}{y_n}$ = $\pm 1$.
Da $y_n \to 0$, haben wir gezeigt, dass $f$ nicht differenzierbar ist. \medskip

\textbf{Ex. Sei $f$ eine Funktion, die in $x_0$ differenzierbar ist. Sei $n \geq 2$, berechne $\lim_{h \to 0} \frac{f(x_0 + nh) - f(x_0 + (n-2)h)}{h}$.} 

Für $n \geq 0$ haben wir $\lim_{h \to 0} \frac{f(x_0 + nh) - f(x_0)}{h} = n \cdot f'(x_0)$. Somit erhalten
wir für $n \geq 2$ $\lim_{h \to 0} \frac{f(x_0 + nh) - f(x_0 + (n-2)h)}{h} = 2 \cdot f'(x_0)$. \medskip

\textbf{Ex. Zeige, dass $f$ gerade $\Rightarrow f'$ ungerade.}

Es gilt zu zeigen, dass $f'(-x) = -f'(x)$. $f'(-x) = \lim_{h \to 0} \frac{f(-x + h) - f(-x)}{h} =
\lim_{h \to 0} \frac{f(x - h) - f(x)}{h} = \lim_{h \to 0} \frac{f(x + h') - f(x)}{-h'} = -f'(x)$\medskip

\textbf{Ex. Zeige, dass die Funktion $f : \R \to \R, x \to e^x-1-x$ nur für $x = 0$ verschwindet.} 

Es gilt offensichtlich $f(0) = 0$. $f'(x) = e^x - 1 > 0$ für $x > 0$. Dass heisst, dass $g$ auf $]0, \infty[$
strikt monoton wächst. Analoges gilt für $x < 0$, deshalb ist $x = 0$ die einzige Nullstelle von $f$.
\medskip

\textbf{Ex. Zeige, dass für $f(x) = \sin(x + a)$ gilt $f^{(n)}(x) = \sin(x + a + \frac{n\pi}{2})$.} 

Wir verwenden Induktion. Für $n = 1$ gilt: $f'(x) = \cos(x + a) = \sin(x + a + \pi/2)$.
Für den Induktionsschritt gilt: $f^{(n)} = (\sin(x + a + \frac{(n-1)\pi}{2}))' =
\cos(x + a + \frac{(n-1)\pi}{2}) = \sin(x + a + \frac{n\pi}{2})$. \medskip

\textbf{Ex. Gegeben die Funktion $f : x \to x^x$ benutze die Taylor Approximation im Punkt $x_0 = 1$
zur dritten Ordnung, um eine Approximation von $(\frac{7}{5})^{\frac{7}{5}}$ anzugeben.} 

Wir berechnen zuerst die ersten drei Ableitungen: $f'(x) = (1 + \ln(x))x^x, f''(x) = (1 + \ln(x))^2x^x + \frac{x^x}{x},
f'''(x) = (1 + \ln(x))^3x^x  + \frac{3(1+\ln(x))x^x}{x} - \frac{x^x}{x^2}$. Nun können wir die Taylor
Approximation $f(x) \approx f(1) + f'(1)(x-1) + \frac{f''(x)}{2}(x-1)^2 + \frac{f'''(x)}{6}(x-1)^3$ 
berechnen. Dies ergibt dann $f(\frac{7}{5}) \approx 1 + (\frac{7}{5} - 1) + (\frac{7}{5} - 1)^2 + \frac{1}{2} (\frac{7}{5} - 1)^3
= \frac{199}{125}$. \medskip

\textbf{Ex. Sei $f(x) = \int_0^\infty e^{-xy}\cos(y)dy$ $x \in ]0, \infty[$, zeige, dass $f$ wohldefiniert ist.}

Wir fixieren ein $x > 0$ und müssen nun zeigen, dass für $h(t) = \int_0^t e^{-xy}\cos(y)dy$ ein 
$\Limit h(n)$ exisitert. Wir machen folgende Abschätzung: $|h(T + s) - h(T)| = \int_T^{T+s} e^{-xy}\cos(y)dy
\leq \int_T^{T+s} e^{-xy}dy = \frac{-(e^{-x(T+s)} - e^{-xT})}{x} \leq \frac{1}{e^{xT}x}$. Es folgt,
dass $(h(n))$ eine Cauchy Folge ist und $y = \Limit h(n)$ existiert. Aus $|h(t) - y| \leq |h(t) - 
h(\lfloor t \rfloor)| + |h(\lfloor t \rfloor) - y| \leq \frac{1}{e^{x \lfloor t \rfloor}x} + 
|h(\lfloor t \rfloor) - y| \to^{t \to \infty} 0$ folgt, dass $\Limit h(n) = y$. \medskip

\textbf{Ex. Integriere $\int_0^1 t^2 \cos(2t)dt$}.

Wir wenden zwei mal partielle Integration an. $\int_0^1 t^2 \cos(2t)dt = \frac{1}{2}\sin(2t)t^2|_0^1
- \int_0^1\sin(2t)tdt = \frac{1}{2}\sin(2)+ \frac{1}{2}(\cos(2t)t)|_0^1-\frac{1}{2}\int_0^1\cos(2t)dt
= \frac{1}{2}\sin(2)+ \frac{1}{2}\cos(2) - \frac{1}{4}\sin(2)$. \medskip

\textbf{Ex. Integriere folgende Funktion $f(x) = \frac{t}{t^3 + t^2 - t -1}$.} 

Zuerst verwenden wir Partialbruchzerlegung um $f(x) = \frac{1}{2(t+1)^2} - \frac{1}{4(t+1)} + \frac{1}{4(t-1)}$
zu erhalten. Nun können wir die einzelnen Brüche leicht integrieren. Somit ist $F(x) = \frac{-1}{2(t+1)}
- \frac{1}{4}(\log(t+1) - \log(t-1))$. \medskip

\textbf{Ex. Integriere folgendes uneigentliches Integral: $\int_1^\infty\frac{1}{x^3}\sqrt{\frac{x}{x+1}}dx$.}

Wir formen zuerst um und machen eine Abschätzung, $\frac{1}{x^3}\sqrt{\frac{x}{x+1}} \leq \frac{1}{x^3}$
da $\frac{1}{x^3}$ bekanntlich konvergent ist und unser Integral monoton wachsen ist, können wir schliessen
dass der Grenzwert für das Integral existiert, d.h. es konvergiert. Wir machen die Substitution $t = \sqrt{\frac{x}{x + 1}}$
und erhalten: $\int_1^a\frac{1}{x^3}\sqrt{\frac{x}{x+1}}dx = 2(-\frac{1}{3}(\frac{a}{a+1})^{-3/2}+\sqrt{\frac{a}{a+1}}+ \frac{2^{3/2}}{3}-\sqrt{2})$.
Für $a \to \infty$ ergibt dass $\frac{2}{3}(2- \sqrt{2})$. \medskip

\textbf{Ex. Integriere folgendes uneigentliches Integral: $\int_{-\infty}^\infty\frac{e^{-1/|x|}}{x^2}dx$.}
 
Da sowohl Nenner als auch Zähler immer positiv sind, konvergiert das Integral, genau dann wenn 
$\int_{0}^\infty\frac{e^{-1/x}}{x^2}dx$ konvergiert. Mit der Substitution $t = \frac{1}{x}$ erhalten wir
$\int_{\epsilon}^A\frac{e^{-1/x}}{x^2}dx = -e^{-\frac{1}{\epsilon}}+e^{-\frac{1}{A}}$. Dieser Ausdruck
strebt nach 1 für $\epsilon \to 0^+, A \to \infty$. Somit konvergiert das Integral und der Grenzwert ist
$2 * \int_{0}^\infty\frac{e^{-1/x}}{x^2}dx = 2$. \medskip

\textbf{Ex. Berechne für alle $m \in \N^*$ den Wert des Integrals $\int_0^{\pi / 2} \cos^m(x)dx$.} 

Für $m = 1$ und $m = 2$ haben wir: $\int_0^{\pi / 2} \cos(x)dx = 1, \int_0^{\pi / 2} \cos^2(x)dx = \frac{\pi}{4}$. 
Für $m \geq 3$ gilt: $\int_0^{\pi / 2} \cos^m(x)dx = (m + 1) \int_0^{\pi / 2} \cos^{(m-2)}(x)\sin^2(x)dx
= (m + 1)(\int_0^{\pi / 2} \cos^{(m-2)}(x)dx - \int_0^{\pi / 2} \cos^m(x)dx) = \frac{m-1}{m} \int_0^{\pi / 2} \cos^{(m-2)}(x)dx$.
\medskip

\textbf{Ex. Zeige, dass für jedes $c \in R$ die Funktion $\exp(-t^2)$ auf $]- \infty, c]$ nach $t$ integrierbar ist.} 

Es genügt zu zeigen, dass $\exp(-t^2)$ auf $]- \infty , -1]$ integrierbar ist, da die Funktion auf 
$]-1, c]$ für all $c > -1$ integrierbar ist. Um dies zu zeigen verwenden wir folgende obere Schranke
der Funktion: $\exp(-t^2) < \exp(t)$. Somit ist $\int_{-\infty}^{-1} \exp(t)dt = \lim_{a \to - \infty} \int_a^{-1}\exp(t)dt
= \lim_{a \to -\infty} (e^{-1} - e^a) = e^{-1}$. \medskip

\textbf{Ex. Sei $f : [0,1] \to \R$ mit $x = 0$ für $x = 0$ oder $x \in \R \backslash \Q$ und $\frac{1}{q}$ für $x = \frac{p}{q}$.
Zeige, dass $f$ integrierbar ist.} 

Wir wollen zeigen, dass sowohl Ober- als auch Untersumme gleich sind. Wir merken, dass jede Partition
$[x_{i-1}, x_i]$ eine irrationale Zahl beinhaltet, solange $x_{i-1} \neq x_i$. Damit ist die Untersumme
immer $0$. Sei $\epsilon > 0, n \in \N$, so dass $\frac{1}{n} < \frac{\epsilon}{2}$ und $B_n$ die Menge aller
koprimen Brüche mit Nenner kleiner gleich $n$. Sei $m = \#B_n$, wähle $k > m$ und $P = \{x_0,...,x_k \} $
eine Partition mit max$_{1 \leq i \leq k}|x_i - x_{i-1}| < \frac{\epsilon}{4m}$, wir notieren $P = P_{\frac{\epsilon}{4m}}$.
Da $B_n$ aus $m$ Elementen besteht, gibt es höchstens $2m$ Intervalle, die $B_n$ schneiden. 
Falls $x \notin B_n$, dann ist $x$ entweder irrational oder ein Bruch $\frac{p}{q}$ mit $q > n$. Somit
gilt $f(x) \leq \frac{1}{n}$. Wir schätzen die Obersumme $S(f,P)$ von oben ab: $S(f,P) = 
\sum_{i : [x_{i-1}, x_i] \cap B_n \neq \emptyset}(x_i - x_{i-1})f_i + \sum_{j : [x_{j-1}, x_j] \cap B_n = \emptyset}(x_j - x_{j-1})f_j$.
Die erste Summe läuft über $\leq 2m$ Indizes und es gilt $f_i \leq 1$. In der zweiten Summe können 
wir $f_i$ mit $\frac{1}{n}$ abschätzen. Wir kombinieren und erhalten $S(f,P) \leq 2m \frac{\epsilon}{4m}\cdot 1 + 1 \cdot \frac{1}{n}
< \frac{\epsilon}{2} + \frac{\epsilon}{2} = \epsilon$. Somit haben wir gezeigt $S(f,P) - s(f,P) < \epsilon$
und $f$ ist integrierbar. \medskip

\vfill

\begin{center}
	\includegraphics[width=0.75\linewidth]{you-can-do-it-meme7.jpg}
\end{center}

